\section{Spectral Adaptation as Choosing a Prior}\label{sec:setup}

\paragraph{Setting.} A signal $x$ lives on a continuous domain
$\Omega\subset\mathbb{R}^d$. We observe $y=Ax+n$ with a known linear operator
$A$ (blur and decimation, sparse sampling, or the footprint integrals of
several shifted frames) and $n\sim\mathcal{N}(0,\sigma_n^2 I)$. A continuous
representation $f_\theta$ is fitted by minimising $\|Af_\theta-y\|^2$
(possibly with regularisation or early stopping) and queried anywhere.

\begin{proposition}[Fourier features are a GP prior; Bochner]\label{prop:bochner}
Let $f(u)=\sum_{j=1}^m w_j\varphi_j(u)$ with
$\varphi_j(u)=\sqrt{2/m}\cos(2\pi b_j^\top u+c_j)$, $b_j\sim p(b)$,
$c_j\sim\mathcal{U}[0,2\pi)$ and $w\sim\mathcal{N}(0,s^2I)$. As
$m\to\infty$, $f$ is a zero-mean stationary Gaussian process with covariance
$k(h)=s^2\,\mathbb{E}_b[\cos(2\pi b^\top h)]$ and spectral density
$S(\nu)=s^2[p(\nu)+p(-\nu)]/2$, equal to $s^2p(\nu)$ for symmetric $p$.
\end{proposition}
Hence a ridge read-out of Fourier features is the GP posterior mean under
prior spectrum $S$; with $p=\mathcal{N}(0,\sigma^2I)$,
$S(\nu)\propto\exp(-\|\nu\|^2/2\sigma^2)$. Tuning $\sigma$ (SuperF) or
choosing frequency parameters from a target spectrum (FreSh, WINNER, SEC)
thus picks a prior spectrum within a restricted family; a GP with a fitted or
known spectrum picks it directly. For an MLP on Fourier features trained by
gradient descent, the lazy-training limit gives kernel regression with a
stationary NTK modulated by $p$ \citep{tancik2020fourier,jacot2018ntk}, under additional architectural and optimisation assumptions. General trained MLPs and grid fields do not thereby have known stationary spectral priors; we compare their behaviour empirically
(\S\ref{sec:exp-ab}).

\begin{proposition}[Linear upper bound]\label{prop:bound}
For zero-mean or correctly centred signal and noise, let $x$ have covariance $C_P$, uncorrelated noise covariance $\sigma_n^2I$, and positive-definite $C_y$. Every linear
estimator $\hat x=Gy$ satisfies
$\mathbb{E}\|x-\hat x\|^2\ge\mathbb{E}\|x-\hat x_{\mathrm{LMMSE}}(P)\|^2$,
with $\hat x_{\mathrm{LMMSE}}(P)=C_PA^\top(AC_PA^\top+\sigma_n^2I)^{-1}y$.
\end{proposition}
This optimum holds in expectation under the specified covariance and noise.
Second moments and zero signal--noise cross-covariance suffice for linear
optimality; joint Gaussianity additionally identifies the conditional mean
with MMSE. A generic spectrum does not define a distribution-free ceiling
on real images, and an image periodogram is an oracle reference rather than
an attained optimality certificate. For any fixed linear $H$ and evaluation
map $W$, with $C_y=AC_PA^\top+C_n$ positive definite,
\begin{align*}
&\mathbb E\|W(x-Hy)\|^2-\mathbb E\|W(x-H_Py)\|^2\\
&=\operatorname{tr}[W(H-H_P)C_y(H-H_P)^\top W^\top]\ge0,
\end{align*}
where $H_P=C_PA^\top C_y^{-1}$. The same covariance-conditioned optimum
minimises total and projected linear risk.

\begin{definition}[Per-band linear recoverability]
For a band projector $\Pi_b$,
\[
\Rlin(A,\sigma_n,P)=1-\frac{\mathbb{E}\|\Pi_b(x-\hat x_{\mathrm{LMMSE}}(P))\|^2}{\mathbb{E}\|\Pi_b x\|^2}.
\]
\end{definition}
$\Rlin$ depends only on $A$, $\sigma_n$ and $P$. With the generic natural-image
prior $P\propto1/f^2$ it requires no ground truth and no training; we compute
it by sampling from the prior, simulating the observation and applying the
LMMSE, with explicit reference-sampling and solve costs. Its amplitude is an additional prior assumption. When a trained or otherwise non-linear estimator
is evaluated on real images, we report the same quantity measured on data,
$R_b=1-\|\Pi_b(x-\hat x)\|^2/\|\Pi_b(x-\bar x)\|^2$ using the colour aggregation specified by each protocol;
for a trained network we write $\Rnet$.

\paragraph{Bands.} Image bands are square annuli in HR pixel units, indexed by $\max(|\nu_x|,|\nu_y|)$,
$b_1=[0,\tfrac18)$, $b_2=[\tfrac18,\tfrac14)$, $b_3=[\tfrac14,\tfrac12]$
cycles/pixel. At $\times4$, $b_1$ lies below the LR Nyquist limit and
$b_2$, $b_3$ lie one and two octaves above it.

\paragraph{Closed form for decimation.} For ideal periodic $A$ = blur $K$ followed by
decimation by $s$, each LR frequency $\omega$ collects $s^d$ HR frequencies
$\nu_m=\omega+m/s$, and
\begin{equation}\label{eq:alias}
R(\nu_0)=\frac{|\hat K(\nu_0)|^2P(\nu_0)}{\sum_m|\hat K(\nu_m)|^2P(\nu_m)+N_\omega}.
\end{equation}
The sum in the denominator is aliasing. In single-image and sparse-pixel
settings it dominates the sensor noise, which is why $\Rlin$ there is nearly
insensitive to $\sigma_n$ (Table~\ref{tab:map}).

Here $N_\omega$ is noise power under a compatible Fourier normalisation.
Changing other powers within the alias set changes the effective noise;
the scalar mismatch identity assumes that effective noise is fixed.
Finite crops, boundary interpolation, rotations and sparse sampling
normally require the matrix risk expression. F1/G1 use HR256 and 73 images;
L1/L2 use HR240 and 74, with their own window and pooling conventions.
Empirical recovery can be negative and need not equal reference recovery.

Zero-noise singular cases require a pseudoinverse or explicitly stated
regularisation. A numerical noise floor changes the estimator; its value
must not replace the physical observation-noise covariance in risk evaluation.
